\documentclass[11pt]{article}
\usepackage[margin=2.2cm]{geometry}
\usepackage{amsmath,amssymb,mathtools}
\usepackage[shortlabels]{enumitem}
\setlist[enumerate,1]{label=\textbf{\arabic*.}, leftmargin=*, itemsep=10pt, topsep=6pt}
\usepackage{fancyhdr}
\usepackage{xcolor}

\newcommand{\R}{\mathbb{R}}
\newcommand{\E}{\mathbb{E}}
\DeclareMathOperator{\Var}{Var}
\DeclareMathOperator{\Cov}{Cov}
\newcommand{\ans}[1]{\par\noindent\textbf{\color{blue!50!black}Answer: #1}}

\pagestyle{fancy}
\fancyhf{}
\fancyhead[L]{Mathematical Finance \& Stochastic Processes}
\fancyhead[R]{Practice Quiz A --- Solutions}
\fancyfoot[C]{\thepage}
\renewcommand{\headrulewidth}{0.4pt}

\begin{document}

\begin{center}
{\LARGE\bfseries Practice Quiz A --- Solutions}
\end{center}
\vspace{6pt}
\hrule
\vspace{10pt}

\begin{enumerate}

\item[\textbf{Q1.}] \ans{A, B, D.}
\begin{itemize}[nosep]
\item[A] $T_1$ is the hitting time of a continuous, adapted process at a fixed level: a standard stopping time.
\item[B] $B_t-(2-3t)=B_t+3t-2$ is continuous and adapted, so $T_2$ is the hitting time of level $0$ by a continuous adapted process -- a stopping time (cf.\ Sheet 3, Problem 3.4(a)).
\item[C] $T_3$ is a \emph{last-exit} time (a supremum, not an infimum); deciding $\{T_3\le t\}$ requires knowing $B$ has no further zero in $(t,1]$, which is not $\mathcal F_t$-measurable. Not a stopping time.
\item[D] The minimum of a stopping time and a deterministic constant is always a stopping time: $\{T_4\wedge5\le t\}=\{T_4\le t\}\cup\{t\ge5\}\in\mathcal F_t$.
\end{itemize}

\item[\textbf{Q2.}] \ans{B.} This mirrors Sheet 3, Problem 3.4(b): the time of the maximum is a last-exit-type random time and is not a stopping time.

\item[\textbf{Q3.}] \ans{C.} By Problem 3.1(b)/4.1(i)-type reasoning, $\exp(\theta B_t-\theta^2t/2)$ is a martingale for every $\theta$. Matching $-3t=-\theta^2t/2$ gives $\theta^2=6$, i.e.\ $\theta=\pm\sqrt6$.

\item[\textbf{Q4.}] \ans{A.} For a continuous Gaussian martingale with $\Gamma(t,t)=\E[M_t^2]=4t$, the process $\exp(\theta M_t-\theta^2\Gamma(t,t)/2)$ is a martingale for every $\theta$ (Problem 4.1(i)). Here $\theta^2\Gamma(t,t)/2=\theta^2\cdot4t/2=2\theta^2t$. Setting $2\theta^2t=8t$ gives $\theta^2=4$, so $\theta=\pm2$.

\item[\textbf{Q5.}] \ans{C.} By Problem 4.1(h), $\langle M,M\rangle_t=\Gamma(t,t)=\E[M_t^2]=4t$.

\item[\textbf{Q6.}] \ans{C.} For $Y_t=\int_0^t H_s\,dB_s$, $\langle Y,Y\rangle_t=\int_0^t H_s^2\,ds$. Here $H_s=s^2$, so $\langle Y,Y\rangle_t=\int_0^t s^4\,ds=t^5/5$.

\item[\textbf{Q7.}] \ans{A, B, D.}
\begin{itemize}[nosep]
\item[A] $B_s$ is $\mathcal F_s$-measurable and $(B_t-B_s)^2$ is independent of $\mathcal F_s$ with mean $t-s$, so $\E[B_s(B_t-B_s)^2]=\E[B_s]\,(t-s)=0$ (Problem 3.4(c)).
\item[B] Writing $B_t=B_s+(B_t-B_s)$: $\E[B_s^2B_t]=\E[B_s^3]+\E[B_s^2]\E[B_t-B_s]=0+s\cdot0=0$ (Problem 3.4(d)). Equivalently: any odd-order moment of jointly Gaussian mean-zero variables vanishes.
\item[C] $\E[B_sB_t]=\Cov(B_s,B_t)=s\wedge t=s\ne0$ in general. \textbf{Not} zero.
\item[D] $B_s\sim N(0,s)$, so all odd moments vanish: $\E[B_s^3]=0$.
\end{itemize}

\item[\textbf{Q8.}] \ans{B.} Since all processes are centred, $\E[(B_t-B_s)B_u]=\Cov(B_t-B_s,B_u)=\Cov(B_t,B_u)-\Cov(B_s,B_u)=(t\wedge u)-(s\wedge u)=t-s$ for $s<t<u$ (Problem 3.4(e)).

\item[\textbf{Q9.}] \ans{C.} $B_s,B_t,B_u$ are jointly Gaussian and mean zero, so the third (odd) moment of their product vanishes: $\E[B_sB_tB_u]=0$ (Problem 3.4(f)).

\item[\textbf{Q10.}] \ans{A.} The increments $B_t-B_s$ and $B_v-B_u$ come from the disjoint intervals $(s,t]$ and $(u,v]$ (since $t<u$), hence are independent, mean-zero random variables; their product has expectation $0$.

\item[\textbf{Q11.}] \ans{A.} $\Cov(B_t,Y_t)=\E[B_tY_t]=\E\Big[\int_0^t1\,dB_s\int_0^t2s\,dB_s\Big]=\int_0^t1\cdot2s\,ds=t^2$ (Itô isometry / covariance formula for Itô integrals driven by the same $B$).

\item[\textbf{Q12.}] \ans{A.} For a continuous Gaussian martingale, $\Cov(M_s,M_u-M_t)=0$ for $s\le t<u$ (Problem 4.1(e)), which implies $\Cov(M_s,M_t)=\Gamma(s,s)=\Gamma(s\wedge t,s\wedge t)$ for $s\le t$. Here $\Cov(M_2,M_7)=\Gamma(2,2)=4\cdot2=8$.

\item[\textbf{Q13.}] \ans{A.} A deterministic, square-integrable integrand produces a genuine (not just local) martingale; $\int_0^\infty e^{-4s}ds<\infty$ guarantees square-integrability on any finite horizon.

\item[\textbf{Q14.}] \ans{B.} Write $X_t=e^{-2t}Z_t$ with $Z_t=\int_0^te^{2s}\,dB_s$ a martingale, $dZ_t=e^{2t}dB_t$. Then by the product rule, $dX_t=-2e^{-2t}Z_t\,dt+e^{-2t}dZ_t=-2X_t\,dt+dB_t$. The nonzero drift term $-2X_t\,dt$ shows $X$ is \emph{not} a martingale (it is the stationary Ornstein--Uhlenbeck process). This illustrates that ``integrand depends on the free variable $t$ itself'' breaks the martingale property, in contrast to Q13 and Problem 5.5(c)(i).

\item[\textbf{Q15.}] \ans{A.} By the Itô isometry, $\Var(X_t)=\E[X_t^2]=\int_0^t(3s+1)^2\,ds=\int_0^t(9s^2+6s+1)\,ds=3t^3+3t^2+t$.

\item[\textbf{Q16.}] \ans{A.} By the Itô isometry, $\E[Z_t^2]=\E\Big[\int_0^tY_s^2\,ds\Big]=\int_0^t\E[Y_s^2]\,ds$. Since $\E[Y_s^2]=\int_0^su^2\,du=s^3/3$ (again Itô isometry), $\E[Z_t^2]=\int_0^t\frac{s^3}{3}\,ds=\frac{t^4}{12}$.

\end{enumerate}

\end{document}
